\chapter{Requirement Analysis}

The objective of the requirement analysis is to establish clear expectations of the PCMA's joint design for the purposes of concept development.

\section{Identification of Stakeholders}

The first step in the requirement analysis process is to identify stakeholders from the operational and support life cycle phases of the project. Stakeholders will be classified into the following categories, as defined below:
\begin{itemize}
	\item Primary stakeholders directly interact with the system while it is in operation.
	\item Secondary stakeholders contribute to the development and functioning of the system.
	\item Indirect stakeholders rely on or are affected by the system, but do not interact with the system directly.
\end{itemize}
The stakeholders relevant to the project are defined in Table~\ref{tab:SID table} below.
\begin{table}[htbp]
	\centering
	\caption{Identified stakeholders and their classifications}
	\label{tab:SID table}
	\begin{tabular}{|l|l|l|}
		\hline
		\textbf{Stakeholder ID} & \textbf{Description} & \textbf{Classification} \\ \hline
		SID\_01 & Supervisor of the project & Primary \\ \hline
		SID\_02 & Students researching PCMA development & Secondary \\ \hline
		SID\_03 & Companies invested in PCMA development & Indirect\\ \hline
		SID\_04 & Mechanical workshop staff & Indirect \\ \hline
		SID\_05 & Mechatronics laboratory managers & Indirect \\ \hline
	\end{tabular}
\end{table}

\section{System Context}

The next step is to define the narrower system of interest (NSoI) and the wider system of interest (WSoI) of the project\footnote{The NSoI is the system being designed and the WSoI includes the NSoI and elements that the NSoI interacts with, \citep{Callister:2011}.}, \citep{Callister:2011} for the purposes of this project, the joint was identified as the NSoI and WSoI is the NSol with a reliable power supply and an artefact with which the joint will be tested. Figure~\ref{fig:system-context} illustrates the system context, with the NSoI shown as a "black box", \citep{Callister:2011}.

\begin{figure}[h]
	\centering
	\includegraphics[width=0.6\linewidth]{"../../../../../../../../../Downloads/System Context"}
	\caption{System Context}
	\label{fig:system-context}
\end{figure}

\section{Design Independent Parameters}

Next, the design independent parameters (DIPs) are defined. DIPs influence the system during the operational and support life cycle stages, but are not controllable nor are part of the system being defined, \citep{Callister:2011}. The DIPs of the joint design are detailed in Table~\ref{tab:DIPs} below.

\begin{table}[h]
	\centering
	\caption{Design Independent Parameters}
	\label{tab:DIPs}
	\begin{tabular}{|l|l|l|}
		\hline
		\textbf{DIP\_ID} & \textbf{Description} & \textbf{Expected range} \\\hline
		DIP\_01.1 & Electrical supply & \SI{12}{V} - \SI{24}{V} \\\hline
		DIP\_01.2 & Electrical supply & \SI{12}{A} maximum \\\hline
		DIP\_01.3 & Electrical supply & Direct current \\\hline
		DIP\_02.1 & Artefact & Compliant with ISO 10360-12:2016 \\\hline
		DIP\_02.2 & Artefact & 20 points reference points \\\hline
		DIP\_03.1 & Environment & Laboratory ambient temperature: \SI{0}{\degree C} - \SI{35}{\degree C} \\\hline
		DIP\_03.2 & Environment & Laboratory ambient humidity: \SI{20}{\%} - \SI{80}{\%} \\\hline
		DIP\_{03.3} & Project budgect & R 6000 \\\hline
	\end{tabular}
\end{table}

\section{Functional Requirements}

In this section, the primary and top-level derived functional requirements (FRs) for the joint design are defined. The only FRs that are considered are only associated with joint's operation. FRs are categorised into primary and derived FRs, in which:
\begin{itemize}
	\item Stakeholders are directly concerned with primary FRs.
	\item Derived FRs are required so that the systen can perform primary FRs.
\end{itemize}
As seen in Table~\ref{tab:FRs} below, the FRs for the joint design are listed.
\begin{table}[h]
	\centering
	\caption{Table 3: Functional Requirements}
	\label{tab:FRs}
	\begin{tabular}{|l|l|l|}
		\hline
		\textbf{FR ID} & \textbf{Description} & \textbf{Primary / Derived} \\\hline
		FR\_01 & Articulate output link & P \\\hline
		FR\_02 & Hold position & P \\\hline
		FR\_03 & Measure angular displacement & P\\\hline
		FR\_04 & Store angular displacement data & D \\\hline
		FR\_05 & Release position & P \\\hline
	\end{tabular}
\end{table}
\noindent In Figure~\ref{fig:functional-decomposition-2}, the functional requirements of the joint are depicted below.

\begin{figure}[h]
	\centering
	\includegraphics[width=0.7\linewidth]{"../../../../../../../../../Downloads/Functional decomposition (2)"}
	\caption{Top Level Functional Decomposition}
	\label{fig:functional-decomposition-2}
\end{figure}

\section{Technical Performance Measures}

This section defines the technical performance measures (TPMs) which will be used to determine how well the joint design (NSoI) performs the required functions\footnote{TPMs are measurable and unambiguous parameters that can be influenced by designing the joint}. The TPMs, listed in Table~\ref{tab:TPMs}, were derived from the results obtained from the generalised error-propagation model. 

\begin{table}[h]
	\centering
	\caption{Technical Performance Measures}
	\label{tab:TPMs}
	\begin{tabular}{|l|l|l|l|}
		\hline
		\textbf{TPM ID} & \textbf{Related FR} & \textbf{Description} & \textbf{Target range} \\\hline
		TPM\_01 & FR\_01, FR\_05 & Operational articulation angle & \SI{90}{\degree} \\\hline
		TPM\_02 & FR\_01, FR\_05 & Run-out tolerance & \SI{1.571e-3}{rad} \\\hline
		TPM\_03 & FR\_02 & Holding torque & \SI{7.725}{Nm} \\\hline
		TPM\_04 & FR\_03, FR\_04 & Encoder accuracy & \SI{1000}{PPR} \\\hline
		TPM\_05 & FR\_01, FR\_02, FR\_05 & Joint life cycle & \num{e6} life cycles \\\hline
	\end{tabular}
\end{table}